\chapter{Análisis y Conclusiones}

Este trabajo aborda el problema de la inferencia de relaciones comerciales entre Sistemas Autónomos (AS). Tarea fundamental para comprender el funcionamiento de Internet, optimizar su desempeño y detectar posibles anomalías. Se construyó un \textit{benchmark} que integra métodos previos (heurísticos y basados en Machine Learning), con el objetivo de establecer una línea base de comparación. Luego, se exploró la resolución de la tarea utilizando Redes Neuronales de Grafos (GNNs), evaluando su desempeño frente al \textit{benchmark}.\\


% RESULTADO BENCHMARK ------------------------------------

En el \textit{benchmark}, los métodos con mejor Accuracy corresponden a los métodos \textbf{ProbLink} (94,83 \%) y \textbf{BGP2Vec} (94,51 \%). Esto nos da a entender que enfoques probabilísticos y de Machine Learning (incluso en dos etapas, como BGP2Vec) capturan de mejor forma las características de los caminos de BGP.
Los métodos de BGP2Vec y Problink toman en cuenta tanto la información local (sus vecinos inmediatos) como la estructura global de la topología para etiquetar las relaciones; en cambio,  en contraste con los métodos de Gao\cite{InferringASRelatioships2001} y  Ruan\cite{Inferring-Multilateral-Peering}, realizan la inferencia basándose únicamente en la estructura observada en cada camino individual. 
Esto nos podría decir que, para inferir correctamente los tipos de relaciones entre Sistemas Autónomos, es fundamental tener el contexto global del enlace dentro de la topología de Internet, no es suficiente con conocer solo su posición dentro de un camino BGP (\texttt{AS\_PATH}) para una mejor inferencia.\\


% GNN (por etapas) -------------------------------------

Al pasar a las Redes Neuronales de Grafos (GNNs), específicamente \textit{por etapas} (donde primero se preentrenan los \textit{embeddings} y luego se utiliza un clasificador) se obtienen buenos resultados, con una Accuracy de alrededor de 85\%, superando a los métodos tradicionales de inferencia. Aquí también podemos notar que el uso de atributos de \textbf{PeeringDB}~\cite{PeeringDB}, en contraste con utilizar únicamente el grado, mejora la Accuracy (de 81.09\% a 85\%), dando a entender que las características asociadas de los ASes aportan información relevante para la inferencia. \\

Sin embargo, al utilizar los \textit{embeddings} generados por \textbf{DeepWalk} y \textbf{BGP2Vec}, que se construyen exclusivamente a partir de la estructura topológica del grafo, se alcanzan los mejores resultados dentro de este enfoque. Esto indica que la posición estructural de los AS dentro de la topología global de Internet es una fuente de información determinante para inferir correctamente las relaciones entre Sistemas Autónomos.\\

Finalmente, el mejor resultado dentro del enfoque por etapas con GNN se obtuvo al preentrenar los embeddings con la tarea auxiliar de predecir el valor de PageRank, lo que sugiere que incorporar nociones de centralidad y relevancia dentro de la topología ayuda a aprender representaciones más informativas para la clasificación, alineándose con lo dicho previamente.\\


% GNN (end-to-end) ------------------------------------
Al evaluar el enfoque \textit{end-to-end}, se obtuvieron los mejores resultados, alcanzando una Accuracy de 93.44\%, comparable con los métodos de mayor Accuracy del \textit{benchmark}.
De manera consistente con el enfoque por etapas, la inclusión de los atributos de PeeringDB~\cite{PeeringDB} resultó beneficiosa, mejorando la Accuracy de 89.01\% a 93.44\%, lo que demuestra que las características asociadas a este dataset enriquecen las representaciones aprendidas por la GNN. 
Por otra parte, en los experimentos diseñados se observó que, al usar atributos aleatorios, el modelo \textit{GCN} mantuvo un rendimiento aceptable (≈ 85\%), mientras que \textit{GraphSAGE} y \textit{GAT} colapsaron, clasificando la mayoría de los enlaces como P2P.\\







Estos resultados respaldan dos conclusiones:
\begin{itemize}
    \item La estructura del grafo por sí sola contiene información suficiente para aproximar la topología y las relaciones entre AS.
    \item La inclusión de atributos como PeeringDB nos permite alcanzar un mejor rendimiento en la inferencia de relaciones comerciales entre Sistemas Autónomos.
\end{itemize}

Además, los resultados obtenidos permiten observar ciertas propiedades estructurales de la topología de Internet y su impacto en el uso de GNNs. Internet muestra una distribución \textit{long-tail} en muchos de sus atributos (grado, centralidad, tráfico), donde pocos nodos (ASes grandes o Tier-1 que conforman el \textit{core}) concentran la mayoría de las conexiones, mientras que la gran mayoría posee una conectividad baja (periferia). Esto provoca que, al entrenar GNNs, los atributos (como el grado) no siempre generen \textit{embeddings} representativos o, en el enfoque \textit{por partes}, se tienda a generalizar. 
Esto explica por qué arquitecturas como GraphSAGE y GAT, que enfatizan la agregación de información desde vecinos seleccionados o ponderados, presentan un rendimiento más inestable, mientras que GCN, al basarse en una agregación más homogénea y en la normalización de la matriz de adyacencia, logra capturar mejor la estructura global incluso en estos escenarios (Caso 3 y 4 del enfoque \textit{por etapas}). \\

Por otro lado, los experimentos también evidencian la ``trampa de la Accuracy'' ya que dado que la clase P2P es la más frecuente en el dataset, un modelo que clasifica todos los enlaces como P2P puede alcanzar una Accuracy alta sin realizar una inferencia real, como en el caso~4 del enfoque \textit{end-to-end}. 
Por ello, métricas como el recall resultan de gran importancia para evaluar esta tarea.\\


Para cerrar, y en vista del trabajo realizado, se propone el siguiente trabajo futuro:
\begin{itemize}
    \item Incorporar atributos dinámicos o temporales, por ejemplo, información proveniente de mensajes \textit{BGP updates}, que permita capturar la evolución temporal de las relaciones entre Sistemas Autónomos.

    \item Ampliar el dataset mediante la integración de nuevas fuentes de datos, como mediciones de \textit{traceroute}, para complementar la visión obtenida desde BGP y construir una representación más completa de la topología de Internet.
    
    \item Evaluar los resultados del modelo en escenarios de detección de anomalías como \textit{route leaks}, \textit{hijacks} u otros eventos anómalos, comparando su rendimiento con métodos existentes.
    
    \item Analizar los \textit{embeddings} generados por las GNN, con el objetivo de estudiar qué propiedades topológicas o atributos capturan y cómo estas representaciones se relacionan con el rol y la posición de los AS dentro del grafo.
    
\end{itemize}
    
























% Primero debemos tener en cuenta que la metrica de Accuracy puede ser engañosa con clases desbalanceadas y en nuestro caso scomo se vio en la Figura X tenmos clases sueprp desbalanceadas, dond ela mayotit a de relaciones son P2Py al rededor d ela mitad p2Cy C2P.
%     Si el 90\% de tus datos son clase "P2P", un modelo que siempre predice "P2P" tendrá 90% de Accuracy… ¡aunque falle completamente las otras clases!


%Para visualizar mejor que esta pasando ocupamos las matrices de confusión, donde podemos ver .... 
% Tus clases son P2P, C2P, P2C. Si una clase es más frecuente, la Accuracy puede esconder errores. Por eso:

%     Usa matriz de confusión para ver qué clases se confunden.

%     Calcula precision, recall y F1-score por clase.

%     La Accuracy puede ser tu métrica global final, pero acompañada de las demás.





% Con estos resultados podemos decir que crear embeddings a partir de link prediction no es una buena metrica para crear la segund atask no rsulva bin y s qud aatascada=o n cirto punto a difrncia d bgp2vc q llga a una mjor acucracy y ovrall rsultados. como sabemos cual sra la mejromtricas para crera mbddinsg?






% CONCLUSION FINAL ---------------------------------------


%El estudio confirma que es factible construir un marco comparativo y competitivo para inferir relaciones AS: los métodos probabilísticos y de DL en dos etapas marcan un listón alto, pero un GNN end-to-end bien diseñado y con atributos adecuados puede situarse muy cerca del estado del arte. La estructura del grafo aporta información sustantiva —en particular con arquitecturas tipo GCN—, aunque la inclusión de atributos de calidad es determinante para superar ambigüedades locales. De cara al futuro, reforzar la capa de datos (más cobertura y temporalidad), enriquecer atributos y adoptar arquitecturas GNN sensibles al tiempo y a la escala de Internet debería traducirse en modelos más estables, precisos y útiles para operadores y analistas de infraestructura.



% RESUMEN -----------------------------------------------
%En síntesis, las contribuciones principales son: (1) la construcción de un benchmark reproducible que integra y estandariza métodos heterogéneos (heurísticos, probabilísticos y de aprendizaje profundo), con evaluación común sobre el mismo subconjunto de CAIDA; (2) la demostración de que un modelo GNN end-to-end, bien alimentado con atributos de PeeringDB, puede aproximar el rendimiento de métodos fuertes como ProbLink y BGP2Vec, manteniendo además interpretabilidad a nivel de señales de grafo; y (3) una caracterización de la sensibilidad a los atributos y a la disponibilidad de paths, útil para futuras campañas de medición.

%\section{Limitaciones y líneas futuras}
%Entre las limitaciones destacan: (i) el uso de RIBs de días específicos/ventanas acotadas, que priorizan el “core” de Internet y pueden sub-representar periferia, afectando la estabilidad temporal; (ii) el desbalance por tipo de relación (p. ej., predominio de c2p/p2c frente a p2p en ciertas vistas), que incide en confusiones observadas en algunos modelos; y (iii) la dependencia de atributos externos (PeeringDB), cuya calidad/actualización impacta directamente el desempeño. Ampliar la recolección (múltiples vantage points y días por mes), incorporar información temporal explícita (series de RIBs) y explorar GNNs temporales o ricas en atención multi-cabeza pueden mejorar robustez y generalización. Integrar señales adicionales (políticas, comunidades BGP, IXP/Route Servers) y técnicas de aprendizaje semi-supervisado/auto-supervisado también parece promisorio para reducir la necesidad de etiquetas y aprovechar la enorme estructura latente del grafo Internet.







%¿Qué pasa con la Accuracy en escenarios desbalanceados?
%La Accuracy mide el total de predicciones correctas sobre todas las clases.

%Si el dataset está desbalanceado (ej: hay muchísimas más relaciones p2p que c2p o p2c), un modelo puede “hacer trampa”: basta con que clasifique casi todo como p2p, y ya tendrá una Accuracy aparentemente alta.

%Pero en realidad estaría ignorando por completo las otras clases, lo que genera un modelo inútil para el problema real.


%Ejemplo ilustrativo:

%Dataset: 80\% p2p, 10\% c2p, 10\% p2c.

%Modelo: predice siempre p2p.

%Accuracy = 80\%.

%Pero recall para c2p y p2c = 0 (nunca las acierta).



%¿Qué hacer para un análisis correcto?

%Tienes que mirar métricas que no dependan tanto del desbalance:

%Recall por clase: te dice si el modelo está detectando las clases minoritarias.

%Precisión por clase: cuántos de los predichos son correctos.

%Matriz de confusión: muestra claramente en qué clases el modelo se equivoca.



%- por q graphs sage tiende a tner mejors resulatdos q las otras? para casos espcificos?
%- por q DP tiend tnmb a tner mejors rsultados q las MLP?

